\section{A Four-Level Model of Semantic Content}
\label{section:4}

%%%%%%%%%%%%%%%%%%%%%%%%%%%%%%%%%%%%%%%%%%%%%%%%%%%%%%%%%%%%%%%%%
% Model Table
%%%%%%%%%%%%%%%%%%%%%%%%%%%%%%%%%%%%%%%%%%%%%%%%%%%%%%%%%%%%%%%%%
% Hierarchy Table
%%%%%%%%%%%%%%%%%%%%%%%%%%%%%%%%%%%%%%%%%%%%%%%%%%%%%%%%%%%%%%%%%
\bgroup
\begin{table*}
\caption{A four-level model of semantic content for accessible visualization.
Levels are defined by the semantic content conveyed by natural language descriptions of visualizations.
Additionally, we offer computational considerations for generating the semantic content at each level of the model.
}

\centering
\def\arraystretch{1.5}

\begin{tabular}{m{0.2cm} m{2.5cm} m{5.5cm} m{6cm}}
\hline
\textsc{\#}
& \textsc{level keywords}
& \textsc{semantic content}
& \textsc{computational considerations} \\
\hline

\cellcolor{Level4}\textcolor{white}{\texttt{4}}
& \cellcolor{Level4Light}\texttt{contextual and domain-specific}
& \cellcolor{Level4Light}{\em domain-specific \mbox{insights}, current events, \mbox{social} and \mbox{political} context, explanations}
& \cellcolor{Level4Light}{contextual knowledge and domain-specific \mbox{expertise} (\emph{perceiver-dependent})}
\\

\cellcolor{Level3}\texttt{3}
& \cellcolor{Level3Light}\texttt{perceptual and cognitive}
& \cellcolor{Level3Light}{\em complex trends, pattern synthesis, \mbox{exceptions}, commonplace concepts}
& \cellcolor{Level3Light}{reference to the rendered visualization and \mbox{``common} knowledge'' (\emph{perceiver-dependent})}
\\

\cellcolor{Level2}\texttt{2}
& \cellcolor{Level2Light}\texttt{statistical and relational}
& \cellcolor{Level2Light}{\em descriptive statistics, extrema, outliers, \mbox{correlations}, point-wise comparisons}
& \cellcolor{Level2Light}{access to the visualization specification or \mbox{backing} dataset (\emph{perceiver-independent})}
\\

\cellcolor{Level1}{\texttt{1}}
& \cellcolor{Level1Light}\texttt{elemental and encoded}
& \cellcolor{Level1Light}{\em chart type, encoding channels, title, axis ranges, labels, colors}
& \cellcolor{Level1Light}{access to the visualization specification or \mbox{rasterized} image (\emph{perceiver-independent})}
\\

\hline
\end{tabular}
\label{table:hierarchy}
\vspace{-5mm}
\end{table*}
\egroup
%%%%%%%%%%%%%%%%%%%%%%%%%%%%%%%%%%%%%%%%%%%%%%%%%%%%%%%%%%%%%%%%%

%%%%%%%%%%%%%%%%%%%%%%%%%%%%%%%%%%%%%%%%%%%%%%%%%%%%%%%%%%%%%%%%%

Our grounded theory process yielded a four-level model of semantic content for the natural language description of visualizations. 
In the following subsections, we introduce the levels of the model and provide example sentences for each.
Table~\ref{table:hierarchy} summarizes the levels, and Table~\ref{table:visualizations} shows example visualizations from our corpus and corresponding descriptions, color coded according to the model's color scale.
Additionally, we offer practical \emph{computational considerations} regarding the feasibility of generating sentences at each level, with reference to the present-day state-of-the-art methods described in Related Work.
While we present them alongside each other for ease of explication, we emphasize that the model levels and computational considerations are theoretically decoupled: the model is indexed to the {semantic content} conveyed by natural language sentences, not to the {computational means} through which those sentences may or may not be generated.

%%%%%%%%%%%%%%%%%%%%%%%%%%%%%%%%%%%%%%%%%%%%%%%%%%%%%%%%%%%%%%%%%
% LEVEL 1
%%%%%%%%%%%%%%%%%%%%%%%%%%%%%%%%%%%%%%%%%%%%%%%%%%%%%%%%%%%%%%%%%
\subsection{Level 1: Elemental and Encoded Properties}
\label{sec:level1}

At the first level, there are sentences whose semantic content refers to elemental and encoded properties of the visualization (i.e., the visual components that comprise a graphical representation's design and construction).
These include the chart type (bar chart, line graph, scatter plot, etc.), its title and legend, its encoding channels, axis labels, and the axis scales. Consider the following sentence (Table~\ref{table:visualizations}.A.1).

\begin{quote}
    \vspace{-1mm}
    \small
    \hlc[Level1]{\textsf{Mortality rate is plotted on the vertical y-axis from 0 to 15\%. Age is plotted on the horizontal x-axis in bins: 10-19, 20-29, 30-39, 40-49, 50-59, 60-69, 70-79, 80+.}}
    \vspace{-1mm}
\end{quote}

\noindent
This sentence ``reads off'' the axis labels and scales as they appear in the bar chart, with no additional synthesizing or interpretation.
Sentences such as this are placed at the lowest level in the model because they refer to content that is \emph{foundational} to visualization construction---comprising the elemental properties of the ``language'' of graphics~\cite{bertin_semiology_1983}.

%----------------------------------------------------------------
% L1 Computation
%----------------------------------------------------------------
\textbf{\emph{Computational Considerations.}}
Semantic content at Level 1 is so foundational that it has long been formalized\,---\,not only theoretically, as in Bertin's \emph{Semiology of Graphics}, but also mathematically and programmatically, as a ``grammar of graphics'' that precisely defines the algorithmic rules for constructing canonical chart types.~\cite{wilkinson_grammar_2005}.
In the case of these construction grammars, Level 1 content is \emph{directly encoded} in the visualization's structured specification (i.e., mappings between data fields and visual properties)~\cite{satyanarayan_vega-lite_2017}.
Thus, for these grammars, generating sentences at Level 1 can amount to ``filling in the blank'' for a pre-defined sentence template. For example, given an appropriate template, the following natural language sentence could be trivially computed using the data encoded in the visualization specification.
\begin{displayquote}
    % \small
    \vspace{-1mm}
    \say{\emph{This is a \texttt{[chart-type]} entitled \texttt{[chart-title]}.
    \texttt{[y-encoding]} is plotted on the vertical y-axis from \texttt{[y-min]} to \texttt{[y-max]}. \texttt{[x-encoding]} is plotted on the horizontal x-axis from \texttt{[x-min]} to \texttt{[x-max]}}.}
    \vspace{-1mm}
\end{displayquote}
\noindent
And similarly, for other sentence templates and elemental properties encoded in the visualization's structured specification.
If the structured specification is not available, however, or if it does not follow a declarative grammar, then \ac{CV} and \ac{NLP} methods have also shown promise when applied to rasterized visualization images (e.g., \textsc{jpeg}s or \textsc{png}s).
For example, recent work has shown that Level 1 semantic content can be feasibly generated provided an appropriate training dataset of pre-defined sentence templates~\cite{kahou_figureqa_2018}, or by extracting a visualization's structured specification from a rasterized visualization image~\cite{poco_reverse-engineering_2017}.

%%%%%%%%%%%%%%%%%%%%%%%%%%%%%%%%%%%%%%%%%%%%%%%%%%%%%%%%%%%%%%%%%
% LEVEL 2
%%%%%%%%%%%%%%%%%%%%%%%%%%%%%%%%%%%%%%%%%%%%%%%%%%%%%%%%%%%%%%%%%
\subsection{Level 2: Statistical Concepts and Relations}

At the second level, there are sentences whose semantic content refers to abstract statistical concepts and relations that are latent the visualization’s backing dataset.
This content conveys computable descriptive statistics (such as mean, standard deviation, extrema, correlations)\,---\,what have sometimes been referred to as ``data facts'' because they are ``objectively'' present within a given dataset~\cite{srinivasan_augmenting_2019, wang_datashot_2020} (as opposed to primarily observed via visualization, which affords more opportunities for subjective interpretation).
In addition to these statistics, Level 2 content includes \emph{relations} between data points (such as ``greater than'' or ``lesser than'' comparisons).
Consider the following sentences (Table~\ref{table:visualizations}.C.2).
\begin{quote}
    \small
    \vspace{-1mm}
    \hlc[Level2]{\textsf{For low income countries, the average life expectancy is 60 years for men and 65 years for women. For high income countries, the average life expectancy is 77 years for men and 82 years for women.}}
    \vspace{-1mm}
\end{quote}

\noindent
These two sentences refer to a statistical property: the computed mean of the life expectancy of a population, faceted by gender and country income-level.
Consider another example (Table~\ref{table:visualizations}.A.2).

\begin{quote}
    \small
    \vspace{-1mm}
     \hlc[Level2]{\textsf{The highest COVID-19 mortality rate is in the 80+ age range, while the lowest mortality rate is in 10-19, 20-29, 30-39, sharing the same rate.}}
    \vspace{-1mm}
\end{quote}

\noindent
Although this sentence is more complex, it nevertheless resides at Level 2. It refers to the \emph{extrema} of the dataset (i.e., the ``highest'' and ``lowest'' mortality rates), and makes two \emph{comparisons} (i.e., a comparison between the extrema, and another between age ranges sharing the lowest mortality rate).
All of the above sentences above share the same characteristic, distinguishing them from those at Level 1: they refer to \emph{relations} between points in the dataset, be they descriptive statistics or point-wise comparisons.
Whereas Level 1 sentences ``read off'' the visualization's elemental properties, Level 2 sentences ``report'' statistical concepts and relations within the chart's backing dataset.

%----------------------------------------------------------------
% L2 Computation
%----------------------------------------------------------------
\textbf{\emph{Computational Considerations.}}
While semantic content at Level 1 requires \emph{only} reference to the visualization's specification, content at Level 2 \emph{also} requires access to the backing dataset.
Here, the two categories of automatic methods begin to diverge in their computational feasibility.
For visualizations with a structured specification, generating sentences at Level 2 is effectively as easy as generating sentences at Level 1: it requires little more computation to calculate and report descriptive statistics when the software has access to the backing dataset (i.e., encoded as part of the visualization specification).
Indeed, many visualization software systems (such as Tableau's Summary Card, Voder~\cite{srinivasan_augmenting_2019}, Quill \textsc{nlg} Plug-In for Power BI, and others) automatically compute summary statistics and present them in natural language captions.
By contrast, for \ac{CV} and \ac{NLP} methods, generating Level 2 sentences from a rasterized image is considerably more difficult\,---\,although not entirely infeasible\,---\,depending on the chart type and complexity.
For example, these methods can sometimes report extrema (e.g., which age ranges exhibit the highest and lowest mortality rates in~\ref{table:visualizations}.A.2)~\cite{obeid_chart--text_2020, demir_generating_2008}.
Nevertheless, precisely reporting descriptive statistics (e.g., the computed mean of points in a scatter plot) is less tractable, without direct access to the chart's backing dataset.


%%%%%%%%%%%%%%%%%%%%%%%%%%%%%%%%%%%%%%%%%%%%%%%%%%%%%%%%%%%%%%%%%
% Visualization & Description Table
%%%%%%%%%%%%%%%%%%%%%%%%%%%%%%%%%%%%%%%%%%%%%%%%%%%%%%%%%%%%%%%%%
% Visualization & Description Table
%%%%%%%%%%%%%%%%%%%%%%%%%%%%%%%%%%%%%%%%%%%%%%%%%%%%%%%%%%%%%%%%%
\bgroup
\begin{table*}
\caption{Example visualizations and descriptions from our corpus. Paragraph breaks in rows \texttt{A} and \texttt{B} indicate a description authored by a unique participant from our corpus gathering survey~\sectionlabel{section:study1design}, while row \texttt{C} shows an curated exemplar description from our evaluation~\sectionlabel{section:study2design}.
% We include the entire corpus, and all of the curated exemplars used in the evaluation, in the supplementary material.
}
\centering
\def\arraystretch{1.5}

\begin{tabular}{m{0.2cm} m{7.5cm} m{9cm}}
\hline
\multicolumn{2}{l}{\textsc{visualization}} & \textsc{description}
\\
\hline

\cellcolor{GrayLight}\textbf{\texttt{A}} &
\includegraphics[scale=0.33]{figures/1134_bar_easy_journalism.png}
\textbf{\texttt{[bar, easy, journalism]}}
&
% \cellcolor{GrayLight}
{\sffamily\smaller

    [1]~\hlc[Level1]{This is a vertical bar chart entitled ``COVID-19 mortality rate by age'' that plots Mortality rate by Age. Mortality rate is plotted on the vertical y-axis from 0 to 15\%. Age is plotted on the horizontal x-axis in bins: 10-19, 20-29, 30-39, 40-49, 50-59, 60-69, 70-79, 80+.}
    [2]~\hlc[Level2]{The highest COVID-19 mortality rate is in the 80+ age range, while the lowest mortality rate is in 10-19, 20-29, 30-39, sharing the same rate.}
    [3]~\hlc[Level3]{COVID-19 mortality rate does not linearly correspond to the demographic age.}

    [4]~\hlc[Level3]{The mortality rate increases with age, especially around 40-49 years and upwards.}
    [5]~\textcolor{white}{\hlc[Level4]{This relates to people's decrease in their immunity and the increase of co-morbidity with age.}}

    [6]~\hlc[Level3]{The mortality rate increases exponentially with older people.}
    [7]~\hlc[Level2]{There is no difference in the mortality rate in the range between the age of 10 and 39.}
    [8]~\hlc[Level3]{The range of ages between 60 and 80+ are more affected by COVID-19.}

    [9]~\textcolor{white}{\hlc[Level4]{We can observe that the mortality rate is higher starting at 50 years old due to many complications prior.}}
    [10]~\hlc[Level3]{As we decrease the age, we also decrease the values in mortality by a lot, almost to none.}
}
\\
\hline

\cellcolor{GrayLight}\textbf{\texttt{B}} &
\includegraphics[scale=0.53]{figures/2220_line_medium_business.png}
\textbf{\texttt{[line, medium, business]}}
& 
% \cellcolor{GrayLight}
{\sffamily\smaller

    [1]~\hlc[Level1]{This is a multi-line chart entitled ``Big Tech Stock Prices'' that plots price by date. The corporations include AAPL (Apple), AMZN (Amazon), GOOG (Google), IBM (IBM), and MSFT (Microsoft). The years are plotted on the horizontal x-axis from 2000 to 2010 with an increment of 2 years. The prices are plotted on the vertical y-axis from 0 to 800 with an increment of 200.}
    [2]~\hlc[Level2]{GOOG has the greatest price over time. MSFT has the lowest price over time.}
    [3]~\hlc[Level3]{Prices of particular Big Tech corporations seem to fluctuate but nevertheless increase over time. Years 2008-2009 are exceptions as we can see an extreme drop in prices of all given corporations.}
    [4]~\textcolor{white}{\hlc[Level4]{The big drop in prices was caused by financial crisis of 2007-2008. The crisis culminated with the bankruptcy of Lehman Brothers on September 15, 2008 and an international banking crisis.}}
    
    [5]~\textcolor{white}{\hlc[Level4]{At the beginning of 2008, every of this stock price went down, likely due to the financial crisis.}}
    [6]~\hlc[Level3]{Then they have risen again and dropped again, more so than previously.}
    
    [7]~\hlc[Level2]{GOOG has the highest price over the years. MSFT has the lowest price over the years.}
    [8]~\hlc[Level3]{GOOG quickly became the richest one of the Big Tech corporations.}
    [9]~\textcolor{white}{\hlc[Level4]{GOOG had experienced some kind of a crisis in 2009, because their prices drop rapidly, but then rebounded.}}
}
\\
\hline

\cellcolor{GrayLight}\textbf{\texttt{C}} &
\includegraphics[scale=0.82]{figures/3310_scatter_hard_academic.png}
\textbf{\texttt{[scatter, hard, academic]}}
&
% \cellcolor{GrayLight}
{\sffamily\smaller

    [1]~\hlc[Level1]{This is a scatter plot entitled ``Born in 2016: Life Expectancy Gap by Gender and Income'' that plots Women Life Expectancy at Birth (Years) by Men Life Expectancy at Birth (Years). The Women Life Expectancy at Birth is plotted on the vertical y-axis from 40 to 90 years. The Men Life Expectancy at Birth is plotted on the horizontal x-axis from 40 to 90 years. High Income Countries are plotted in dark green. Low Income Countries are plotted in light green. A 45 degree line from the origin represents Equal Life Expectancy.}
    [2]~\hlc[Level2]{For low income countries, the average life expectancy is 60 years for men and 65 years for women. For high income countries, the average life expectancy is 77 years for men and 82 years for women.}
    [3]~\hlc[Level3]{Overall, women have a slightly higher life expectancy than men. Women live around 5 to 10 years longer than men. The low income countries are more scattered than the high income countries. There is a visible gap between high and low income countries, indicated by the Income-Age Divide line.}
    [4]~\textcolor{white}{\hlc[Level4]{People living in low-income countries tend to have a lower life expectancy than the people living in high-income countries, likely due to many societal factors, including access to healthcare, food, other resources, and overall quality of life. People who live in lower income countries are more likely to experience deprivation and poverty, which can cause related health problems.}}
}
\\
\hline
\end{tabular}
\label{table:visualizations}
\end{table*}
%%%%%%%%%%%%%%%%%%%%%%%%%%%%%%%%%%%%%%%%%%%%%%%%%%%%%%%%%%%%%%%%%

%%%%%%%%%%%%%%%%%%%%%%%%%%%%%%%%%%%%%%%%%%%%%%%%%%%%%%%%%%%%%%%%%

%%%%%%%%%%%%%%%%%%%%%%%%%%%%%%%%%%%%%%%%%%%%%%%%%%%%%%%%%%%%%%%%%
% LEVEL 3
%%%%%%%%%%%%%%%%%%%%%%%%%%%%%%%%%%%%%%%%%%%%%%%%%%%%%%%%%%%%%%%%%
\subsection{Level 3: Perceptual and Cognitive Phenomena}

At the third level, there are sentences whose semantic content refers to perceptual and cognitive phenomena appearing in the visual representation of the data.
When compared to, and defended against, other forms of data analysis (e.g., purely mathematical or statistical methods), visualization is often argued to confer some unique benefit to human readers.
That is, visualizations do not only ``report'' descriptive statistics of the data (as in Level 2), they also \emph{show} their readers something \emph{more}: they surface unforeseen trends, convey complex multi-faceted patterns, and identify noteworthy exceptions that aren't readily apparent through non-visual methods of analysis (cf., Anscombe's Quartet or the Datasaurus Dozen~\cite{matejka_same_2017}).
Level 3 sentences are comprised of content that refers to these perceptual and cognitive phenomena, usually articulated in ``natural''-sounding (rather than templatized) language.
Consider the following examples (Table~\ref{table:visualizations}.B.3 and \ref{table:visualizations}.C.3, respectively).
\begin{quote}
    \small
    \vspace{-1mm}
    \hlc[Level3]{\textsf{Prices of particular Big Tech corporations seem to fluctuate but nevertheless increase over time. Years 2008-2009 are exceptions as we can see an extreme drop in prices of all given corporations.}}
    \vspace{-1mm}
\end{quote}

\begin{quote}
    \small
    \hlc[Level3]{\textsf{The low income countries are more scattered than the high income countries. There is a visible gap between high and low income countries, indicated by the Income-Age Divide line.}}
    \vspace{-1mm}
\end{quote}

\noindent
These sentences convey the ``overall gist'' of complex trends and patterns (e.g., stock prices ``seem to fluctuate but nevertheless increase''), synthesize multiple trends to identify exceptions (e.g., ``years 2008-2009 are exceptions as we can see an extreme drop'' of multiple graphed lines at that point in time), and do so in ``natural''-sounding language, by referencing commonplace concepts (such as ``fluctuate'', ``extreme drop'', ``visible gap'').
N.b., ``natural''-sounding articulation is necessary but insufficient for Level 3 membership, as it is also possible to articulate Level 1 or 2 content in a non-templatized fashion~\sectionlabel{section:study1results}.

%----------------------------------------------------------------
% L3 Computation
%----------------------------------------------------------------
\textbf{\emph{Computational Considerations.}}
At Level 3, we begin to reach and exceed the limits of present-day state-of-the-art automatic methods.
While there exist ``off-the-shelf'' statistical packages for computing basic trends and predictions in a dataset (e.g., correlations, polynomial regressions, statistical inferences), visualizations allow us to perceive and articulate complex trends for which there may exist no line of ``best fit''.
While automatic methods may eventually approach (or exceed) human capabilities on well-defined tasks~\cite{obeid_chart--text_2020}, for now Level 3 semantic content is likely generated via human (rather than machine) perception and cognition~\cite{morash_guiding_2015}.
Taking inspiration from the \say{mind-independent} versus \say{mind-dependent} ontological distinction~\cite{ali_mind-dependent_2016}, we define sentences at Levels 1 and 2 as \emph{perceiver-independent} (i.e., their content can be generated independently of human or machine perception, without reference to the visualization), while sentences at Level 3 are \emph{perceiver-dependent} (i.e., their content requires a perceiver of some sort; likely a human, although machine perception may increasingly suffice for generating Level 3 content).
Table~\ref{table:hierarchy} summarizes this distinction.

%%%%%%%%%%%%%%%%%%%%%%%%%%%%%%%%%%%%%%%%%%%%%%%%%%%%%%%%%%%%%%%%%
% LEVEL 4
%%%%%%%%%%%%%%%%%%%%%%%%%%%%%%%%%%%%%%%%%%%%%%%%%%%%%%%%%%%%%%%%%
\subsection{Level 4: Contextual and Domain-Specific Insights}

Finally, at the fourth level, there are sentences whose semantic content refers to contextual and domain-specific knowledge and experience.
Consider the following two examples (Table~\ref{table:visualizations}.B.4 and~\ref{table:visualizations}.C.4).
\begin{quote}
    \small
    \vspace{-1mm}
    \textcolor{white}{\hlc[Level4]{\textsf{The big drop in prices was caused by financial crisis of 2007-2008. The crisis culminated with the bankruptcy of Lehman Brothers on September 15, 2008 and an international banking crisis.}}}
    \vspace{-1mm}
\end{quote}

\begin{quote}
    \small
    \textcolor{white}{\hlc[Level4]{\textsf{People living in low-income countries tend to have a lower life expectancy than the people living in high-income countries, likely due to many societal factors, including access to healthcare, food, other resources, and overall quality of life.}}}
    \vspace{-1mm}
\end{quote}

\noindent
These sentences convey social and political explanations for an observed trend that depends on an individual reader's subjective knowledge about particular world events: the 2008 financial crisis and global socio-economic trends, respectively.
This semantic content is characteristic of what is often referred to as ``insight'' in visualization research.
Although lacking a precise and agreed-upon definition~\cite{law_what_2020, tang_extracting_2017, north_toward_2006, chang_defining_2009, law_characterizing_2020}, an insight is often an observation about the data that is complex, deep, qualitative, unexpected, and relevant~\cite{yi_understanding_2008}.
Critically, insights depend on individual perceivers, their subjective knowledge, and domain-expertise.
Level 4 is where the breadth of an individual reader’s knowledge and experience is brought to bear in articulating something ``insightful'' about the visualized data.

%----------------------------------------------------------------
% L4 Computation
%----------------------------------------------------------------
\textbf{\emph{Computational Considerations.}}
{As with Levels 3, we say that Level 4 semantic content is perceiver-dependent, but in a stronger sense}. 
This is because (setting aside consideration of hypothetical future ``artificial general intelligence'') generating Level 4 semantic content is at-present a uniquely human endeavor.
Doing so involves synthesizing background knowledge about the world (such as geographic, cultural, and political relationships between countries), contextual knowledge about current events (e.g., the fact that there was a global recession in 2008), and domain-specific knowledge (e.g., expertise in a particular field of research or scholarship).
However, bespoke systems for narrowly-scoped domains (e.g., those auto-generating stock chart annotations using a corpus of human-authored news articles~\cite{hullman_contextifier_2013}) suggest that some Level 4 content might be feasibly generated sooner rather than later.

Lastly, we briefly note that data-driven predictions can belong to either Level 2, 3, or 4, depending on the semantic content contained therein.
For example: a point-wise prediction at Level 2 (e.g., computing a stock's future expected price using the backing dataset); a prediction about future overall trends at Level 3 (e.g., observing that a steadily increasing stock price will likely continue to rise); a prediction involving contextual or domain-specific knowledge at Level 4 (e.g., the outcome of an election using a variety of poll data, social indicators, and political intuition).
